\section{Discussion}
\label{sec:discussion}

\subsection{Why the Gamma-noise model performed better}

The prescribed epidemic contains two transmission regimes, but the
constant-\(B\) model must explain both with one compromise value.  That
restriction produces a fitted rate that is too low before the change and much
too high afterward.  The Gamma-noise model can instead shift probability
toward lower latent \(B(t)\) paths as post-change reports arrive.  This
adaptation explains both its lower overall RMSE and its much smaller
post-change signed error.

The Gamma transition does not encode the true week-5 boundary or a
deterministic downward trend.  Its conditional mean equals the current rate,
so the fitted decrease emerges through conditioning on the observations:
particle paths that better explain the later reports receive more weight.
The comparison therefore shows the value of allowing transmission to change;
it does not show that a Gamma process generated the prescribed discontinuity.

\subsection{What the result means}

The lower RMSE in 195 of 200 paired data sets shows that a positive latent-rate
model usually recovered this imposed change more accurately than a static-rate
model when both were fitted to the same noisy case series.  Because all three POMPs
retain the same stochastic SIR and observation components, the contrast is
attributable to how the fitting models represent transmission rather than to
one model being stochastic and the other deterministic.

The secondary summaries support the same interpretation while measuring
different features of error.  Absolute task-level mean error can benefit from
cancellation between early negative and later positive errors, so it need not
rank models identically to RMSE in every data set.  The post-fitting likelihood
pattern shows that the more flexible model usually assigned higher likelihood
to the fitted observations, but it does not account for model complexity or
out-of-sample performance.

\subsection{Limitations and next step}

Three limitations define the scope of the conclusion:
\begin{itemize}
\item \emph{Benchmark scope.}  The experiment studies one abrupt path over ten
weeks and conditions on outbreaks satisfying \(\max(H)>20\).  Other change
sizes, timings, smooth trajectories, or weaker outbreaks may yield different
recovery behavior.
\item \emph{Fixed nuisance parameters.}  The recovery rate, reporting model,
population size, and initial state are fixed at their generating values.
Applied analyses must confront confounding between these quantities and
transmission, especially when a flexible \(\sigB\) can absorb other
misspecification.
\item \emph{Finite optimization.}  IF2 uses a finite particle population,
cooling schedule, and set of starting points.  The Gamma model received nine
starts and the constant model six, so the experiment does not establish equal
computational budgets or guarantee global maximization.
\item \emph{Single sampled trajectory.}  Recovery metrics use one
ancestry-preserving latent path from the final filter for each Gamma fit.  They
therefore include finite-particle and path-sampling variation and should not be
interpreted as posterior-mean error or as a summary of trajectory uncertainty.
Repeated trajectory extraction would be needed to quantify that additional
uncertainty.
\end{itemize}

A useful next comparison would replace local Gamma innovations with a smooth
basis representation,
\[
\log B(t)=\sum_{q=1}^{Q}c_q\phi_q(t),
\]
where \(\phi_q(t)\) are B-spline basis functions
\citep{bouman2024timevarying}.  Such a model would trade local stochastic
adaptation for a globally smooth trajectory whose flexibility depends on the
number and placement of knots and on regularization.  Because smoothing can
blur an abrupt change, the basis dimension and penalty would need to be
prespecified and evaluated in a new simulation study.  The next scientific
question is therefore whether Gamma noise or a controlled smooth-basis model
recovers a wider range of plausible transmission paths more reliably.
